\section{Inference and evaluation}
\label{sec:method}

\subsection{Likelihood and maximum-likelihood estimation}

The observed case series does not reveal the epidemic states or, in the
Gamma-noise model, the transmission-rate path.  The observed-data likelihood
therefore averages over all latent state histories compatible with a fixed
parameter vector.  That integral is not available analytically for these
nonlinear stochastic POMP models.  The inferential target is the maximum
likelihood estimate
\begin{equation}
\widehat{\theta}
=
\arg\max_{\theta}\ell(\theta;y_{1:N}).
\label{eq:mle}
\end{equation}
The model-specific static transmission parameters are
\[
\theta_G=(B_0,\sigB),
\qquad
\theta_C=(\Bconst).
\]
The nuisance parameters \(\mu_{IR},N,\rho,k\) and the initial epidemic state
are fixed in this benchmark.  In particular, the latent sequence \(B(t)\) is
not optimized as an unrelated parameter at every time point; it is integrated
over according to the Gamma transition in Equation~\ref{eq:gamma-update}.

\subsection{Particle filtering}

For a fixed parameter vector, a particle filter propagates many possible
latent states through the process model, weights them according to each new
reported count, and resamples states that better explain the observations.
These weighted particles approximate both the observed-data likelihood and
the filtering distribution of the latent state given reports through the
current observation time \citep{gordon1993particle,king2016pomp}.  The analysis
uses the implementation in the \texttt{pomp} package
\citep{king2025pomp64}.

At an observation time, weighting favors particles whose accumulated
incidence makes the reported count plausible; resampling then concentrates
subsequent computation on those histories without collapsing the scientific
model to a single path.  Because only finitely many particles are used, two
independent filters at the same parameter vector need not return identical
likelihood estimates or sampled latent histories.  This particle-filter Monte Carlo
variation is reduced by increasing the particle population and is separate
from every biological or reporting source of randomness in Section~\ref{sec:model}.

\subsection{IF2 estimation}

Iterated filtering (IF2) connects particle filtering to the optimization in
Equation~\ref{eq:mle} \citep{ionides2015if2}.  IF2 repeatedly applies a particle
filter to an extended parameter--state system.  Artificial perturbations allow
the static-parameter particles to explore the likelihood surface, while
observation weighting and resampling favor parameter--state combinations that
better explain the data.  A cooling schedule reduces the perturbation scale
across iterations, producing an approximate maximum likelihood estimate.
Multiple starting points broaden the search but do not guarantee that the
global maximum has been found.

In the implemented search, each starting point underwent a fixed sequence of
600 IF2 iterations.  The decreasing perturbations turn the early iterations
into relatively broad exploration and the later iterations into a more local
search.  Candidate endpoints were not ranked using the final IF2 trace value
alone: five independent ordinary particle filters reevaluated each endpoint
after perturbations had been removed.  If \(\ell_{dsj}\) is the log-likelihood
estimate for data set \(d\), start \(s\), and evaluation \(j\), the ranking
summary was
\begin{equation}
\widetilde\ell_{ds}
=
\log\left\{\frac{1}{5}\sum_{j=1}^{5}\exp(\ell_{dsj})\right\},
\label{eq:logmeanexp}
\end{equation}
computed numerically with \texttt{pomp::logmeanexp}.  Thus the averaging occurs
on the likelihood scale.  It is neither an arithmetic mean of log likelihoods
nor a mean of the latent transmission trajectories.  The endpoint with the
largest finite \(\widetilde\ell_{ds}\) was selected as the best fit.  This
separates parameter search from the final comparison as far as the finite
computation allows.

IF2 perturbations are algorithmic and disappear from the focal model after
optimization.  They are distinct from Gamma process noise, which governs the
scientific evolution of latent \(B(t)\) at the fitted value of \(\sigB\).
Appendix~\ref{app:implementation} records the starting values, particle counts,
perturbation scales, cooling, and selection details.

\subsection{Transmission-rate recovery}

For paired data set \(d\), the selected Gamma parameters were passed to one
final 50,000-particle filter with trajectory ancestry retained.  The
\texttt{filter\_traj()} function then extracted one internally consistent
latent history,
\begin{equation}
\widehat B^{(G)}_d(t_{0:N})
=B^{*}_d(t_{0:N};\widehat\theta_{G,d}),
\label{eq:bhat}
\end{equation}
where the superscript \(*\) denotes the sampled history rather than a posterior
or filtering mean.  Because its ancestry is traced through the final particle
population, this is a finite-particle, plug-in-parameter approximation to a
smoothing trajectory conditional on the complete reported series.  It is not
an exact posterior draw, an uncertainty interval, or a separately optimized
value of \(B\) at each time.  The value at \(t_0\) is retained only for display;
all recovery metrics use the 70 observation times \(t_1,\ldots,t_N\).

The constant-model estimator is the selected static value repeated over those
same observation times,
\[
\widehat B^{(C)}_d(t_n)=\widehat B_{\mathrm{const},d}.
\]
Both models are therefore represented by one observation-time trajectory
before comparison, even though the constant trajectory necessarily repeats a
single value.  No filtering mean enters the Experiment~4 recovery metrics.
For model \(m\in\{G,C\}\), the primary task-level outcome is
\begin{equation}
\operatorname{RMSE}^{(m)}_d
=
\sqrt{
\frac{1}{N_{\mathrm{obs}}}
\sum_{n=1}^{N_{\mathrm{obs}}}
\left[
\widehat B^{(m)}_d(t_n)-\Btrue(t_n)
\right]^2
}.
\label{eq:rmse}
\end{equation}
This definition compares like observation times within each paired data set.
Each data set contributes one RMSE per model and therefore one within-pair
difference.  Averaging these task-level RMSE values gives equal weight to the
200 accepted epidemics, rather than allowing a data set with more cases to
dominate the summary.  Squaring before averaging also prevents early
underestimation and late overestimation from canceling.

Secondary summaries describe signed recovery error.  The task-level mean error
averages \(\widehat B^{(m)}_d(t_n)-\Btrue(t_n)\) over all observation times;
its absolute value measures overall displacement without allowing the final
sign to determine magnitude.  Segment-specific signed errors average the same
quantity before and after week~5.  Post-fitting particle-filter likelihoods
provide an additional same-data fitting diagnostic, but they are neither
penalized nor evaluated on held-out observations and are not treated as a
formal model-selection test.
